\chapter{Conclusiones}\label{ch:conclusions}

El aprendizaje profundo geométrico suele presentarse como un catálogo de arquitecturas: redes convolucionales para imágenes, redes de grafos para moléculas, transformers para secuencias, redes equivariantes para nubes de puntos. Este libro ha seguido el camino contrario. Ha partido de las matemáticas (grupos y sus representaciones, variedades y su curvatura, grafos y sus espectros, geodésicas, fibrados y conexiones) y ha derivado de ellas las arquitecturas, demostrando por el camino los resultados que las justifican. Este capítulo de cierre reúne lo que ese camino ha mostrado, dice con honestidad lo que no ha resuelto y mira hacia dónde va el campo.


\section{La Geometría como Principio Organizador}

\subsection{De las Neuronas a las Simetrías}

El punto de partida fue deliberadamente clásico. El \autoref{ch:neural_networks} formalizó la neurona artificial (\autoref{def:neurona_artificial}), construyó el aparato de teoría de la medida y análisis funcional necesario para demostrar el teorema de aproximación universal (\autoref{sec:aproximacion_universal}) y analizó el descenso de gradiente y su convergencia. Ese capítulo también dejó a la vista la debilidad que motiva todo lo que sigue: una red completamente conectada no tiene forma de saber que una imagen trasladada es la misma imagen, ni que un grafo reetiquetado es el mismo grafo. Tiene que aprender cada simetría a partir de los datos, y lo paga en parámetros y en muestras.

El \autoref{ch:gdl} aportó la estructura que faltaba a través de las cinco Gs (\autoref{subsec:cinco_gs}). La \emph{geometría} fija el dominio sobre el que viven los datos, desde las rejillas euclidianas hasta las variedades diferenciables y la hipótesis de la variedad que enlaza ambas. Los \emph{grupos} describen las simetrías de ese dominio; el teorema de Peter-Weyl (\autoref{thm:peter_weyl}) y el lema de Schur (\autoref{lem:schur}) descomponen sus representaciones en piezas irreducibles, y la equivarianza (\autoref{def:equivariance_gdl}) y la invarianza (\autoref{def:invariance_gdl}) convierten la simetría en una restricción de diseño sobre la red. Los \emph{grafos} discretizan el dominio, y el laplaciano del grafo traslada la teoría espectral del dominio continuo al discreto. Las \emph{geodésicas} proporcionan la noción de distancia a lo largo de la cual se propaga la información, junto con el mapa exponencial y el transporte paralelo. Los \emph{gauges}, mediante fibrados principales (\autoref{def:fibrado_principal_formal}) y conexiones, extienden la equivarianza de una simetría global a una simetría que puede cambiar de un punto a otro.

\subsection{Tres Elecciones y un Único Operador}

Vistas a través de las cinco Gs, las diferencias entre las arquitecturas clásicas se reducen a tres elecciones: el dominio geométrico, el grupo de simetría y la estructura de conectividad. Una vez fijadas, un único operador, el \gls{message_passing} equivariante (\autoref{def:mpnn_unificador}), reproduce cada una de ellas:

\begin{itemize}
	\item Las redes neuronales tradicionales son \glspl{MPNN} sobre grafos completos (\autoref{prop:ann_mpnn}), donde cada neurona intercambia información con todas las demás.
	\item Las \glspl{CNN} realizan \gls{message_passing} sobre rejillas regulares (\autoref{thm:equivalencia_constructiva_cnn_gnn}), con pesos compartidos determinados por el grupo de traslaciones.
	\item Los \glspl{Transformer} implementan \gls{message_passing} adaptativo (\autoref{def:self_attention_basica}), donde los pesos de comunicación se calculan dinámicamente a partir de los datos.
\end{itemize}

Esta unificación no es meramente conceptual. El test de Weisfeiler-Leman (\autoref{thm:wl_mpnn_equivalencia}) establece límites expresivos precisos para las \glspl{MPNN}, mostrando que su poder de distinción de grafos coincide exactamente con el del algoritmo combinatorio 1-WL. La jerarquía de generalización (\autoref{rem:jerarquia_generalizacion}) formaliza las relaciones de inclusión entre las distintas familias arquitectónicas, proporcionando un mapa de sus capacidades relativas, y la síntesis de la \autoref{sec:sintesis_marco_unificador} muestra que el mapa es exhaustivo: cada arquitectura tratada en el libro ocupa una posición en él.

\subsection{Convolución y Atención más allá del Espacio Euclidiano}

Las dos familias que dominan la práctica se reconstruyeron después sobre dominios no euclidianos. El \autoref{ch:cnn} construyó una progresión de generalizaciones convolucionales de complejidad creciente: convoluciones de grupo, redes steerable, convoluciones esféricas, redes $\text{E}(3)$/$SE(3)$-equivariantes, redes gauge equivariantes y convoluciones Clifford-valuadas. Esta cadena culmina en el teorema de universalidad equivariante (\autoref{theorem:equivariant_universality_cnns}), que garantiza que las \glspl{CNN} equivariantes pueden aproximar cualquier función continua equivariante, extendiendo así el resultado clásico de aproximación universal al contexto geométrico. Las extensiones topológicas, desarrolladas mediante los laplacianos de Hodge (\autoref{def:hodge_laplacians_cnn}), proporcionaron herramientas para operar sobre estructuras de orden superior (complejos simpliciales y celulares).

El \autoref{ch:attention} hizo lo mismo con la atención. Su formalización como operador integral la conectó con el análisis funcional del \autoref{ch:neural_networks}; la atención equivariante mediante álgebra de Clifford (\autoref{sec:clifford_attention}) opera directamente sobre multivectores, preservando las simetrías del álgebra geométrica; la atención hiperbólica (\autoref{def:atencion_hiperbolica}) extiende el mecanismo a espacios de curvatura negativa constante, adecuados para datos jerárquicos; y la atención simplicial y de Hodge lo lleva a estructuras de orden superior. El análisis armónico actuó como herramienta unificadora a lo largo de ambos capítulos, conectando las perspectivas espacial y espectral de cada operación geométrica.

\subsection{La Teoría se Encuentra con las Arquitecturas}

Nada de esto se quedó en abstracto. AlphaFold~\citep{jumper2021alphafold} ejemplifica la equivarianza $SE(3)$ en predicción de estructura de proteínas. DeepSets y PointNet ilustran la invarianza a permutaciones, con el teorema de universalidad correspondiente (\autoref{thm:deepsets_universal}) garantizando su capacidad expresiva sobre funciones de conjuntos. Los \gls{GATr} demuestran la aplicabilidad de las álgebras de Clifford en arquitecturas de atención. La tabla de equivalencias \gls{CNN}-atención (\autoref{tab:cnn_attention_equivalencias_fundamentales}) sintetiza las correspondencias formales entre ambos paradigmas, mostrando que comparten estructura matemática a pesar de sus diferencias aparentes, y el teorema de correspondencia (\autoref{thm:correspondencia_cnn_attention}) junto con la conjetura de convergencia (\autoref{conj:convergencia_cnn_attention}) formalizan la relación: bajo condiciones apropiadas, convolución y atención convergen al mismo operador.


\section{Limitaciones y Problemas Abiertos}

\subsection{Poder Expresivo y la Barrera de Weisfeiler-Leman}

La equivalencia entre \glspl{MPNN} y el test 1-WL (\autoref{thm:wl_mpnn_equivalencia}) establece una barrera expresiva fundamental: las \glspl{MPNN} estándar no pueden distinguir grafos que el test 1-WL no distingue. Las extensiones $k$-WL superan esta limitación, pero con un coste computacional que crece exponencialmente con $k$, lo que las hace impracticables para grafos de tamaño moderado. Un problema abierto central es el diseño de arquitecturas con coste polinomial que superen la barrera 1-WL, posiblemente mediante el uso de features estructurales adicionales o mecanismos de atención de orden superior.

La propia barrera se ha refinado en los últimos años. Para grafos discretos, \citet{zhang2023biconnectivity} demuestran que ninguna arquitectura de la clase 1-WL puede decidir la biconectividad (vértices y aristas de corte), una propiedad que un esquema de paso de mensajes sobre distancias generalizadas recupera con coste polinomial, y \citet{zhang2024beyondwl} sustituyen la comparación binaria con la jerarquía $k$-WL por la \textit{expresividad por homomorfismos}, que caracteriza exactamente qué recuentos de subestructuras puede calcular cada familia arquitectónica; para grafos geométricos, el test WL geométrico de \citet{joshi2023expressive} (\autoref{subsubsec:wl_geometrico}) muestra que las capas invariantes son estrictamente menos expresivas que las equivariantes y que el poder expresivo depende de la profundidad, del orden tensorial y del orden de cuerpo de los mensajes. El \gls{oversquashing}, la obstrucción complementaria, se entiende ahora de forma cuantitativa: \citet{digiovanni2023oversquashing} acotan la sensibilidad de la representación de un nodo a las features de entrada de un nodo lejano mediante una cantidad gobernada por la resistencia efectiva (equivalentemente, el tiempo de conmutación) entre ambos nodos, de modo que es la topología del grafo, y no la profundidad o la anchura de la red, la que determina qué pares de nodos pueden intercambiar información.

\subsection{Tensión entre Rigidez y Adaptabilidad}

El trade-off entre rigidez y adaptabilidad (\autoref{rem:tradeoff_rigidez_adaptabilidad}), formalizado en el \autoref{ch:attention}, identifica una tensión fundamental en el diseño de arquitecturas geométricas. Las \glspl{CNN} equivariantes incorporan simetrías de forma rígida, lo que proporciona eficiencia en datos y garantías teóricas, pero limita su capacidad de adaptación cuando las simetrías son solo aproximadas. Los mecanismos de atención ofrecen flexibilidad para aprender relaciones dependientes del contexto, pero a costa de mayor demanda de datos y menor regularización estructural. La conjetura de convergencia \gls{CNN}-atención (\autoref{conj:convergencia_cnn_attention}) sugiere que ambos paradigmas convergen bajo condiciones apropiadas, pero el punto óptimo en el espectro rigidez-adaptabilidad para cada dominio concreto permanece como un problema abierto.

Entre ambos polos ha surgido una tercera opción, que desacopla la equivarianza de la arquitectura. El \textit{frame averaging} \citep{puny2022frameaveraging} convierte una red arbitraria $f$ en exactamente equivariante promediándola sobre un pequeño conjunto de elementos del grupo dependiente de la entrada $F(\mathbf{x}) \subset G$ (un \textit{frame}, que satisface $F(g\mathbf{x}) = g\,F(\mathbf{x})$), $\Phi(\mathbf{x}) = \frac{1}{|F(\mathbf{x})|}\sum_{g \in F(\mathbf{x})} \rho_{\mathcal{Y}}(g)\, f(g^{-1}\mathbf{x})$, sin restringir las capas de $f$ y sin perder universalidad; la canonicalización aprendida \citep{kaba2023canonicalization} reduce el frame a un único elemento, predicho por una red auxiliar equivariante que lleva la entrada a una pose canónica (véase la \autoref{subsec:canonicalizacion_frames}). Cuando el propio grupo es desconocido, los métodos de descubrimiento de simetrías como LieGAN \citep{yang2023liegan} aprenden los generadores de un álgebra de Lie de simetrías a partir de los datos mediante entrenamiento adversario, aportando el ingrediente que necesitaría una cadena completamente automática de los datos al grupo de simetría y de este a la arquitectura (\autoref{subsec:hacia_teoria_unificada}).

\subsection{Gauge Equivarianza y Topología Global}

La teoría gauge desarrollada en el \autoref{ch:cnn} mediante fibrados principales (\autoref{def:fibrado_principal_formal}) proporciona un marco riguroso para las simetrías locales. Sin embargo, las obstrucciones topológicas globales (fibrados no triviales, holonomía no trivial alrededor de ciclos) plantean dificultades tanto teóricas como computacionales. La brecha entre la elegancia del formalismo gauge y la implementación de arquitecturas eficientes a escala constituye un desafío práctico significativo: las \glspl{CNN} gauge equivariantes actuales requieren discretizaciones cuidadosas que pueden no preservar las propiedades topológicas globales del espacio subyacente.

\subsection{Estructuras de Orden Superior}

Las herramientas topológicas desarrolladas en este libro (laplacianos de Hodge (\autoref{def:hodge_laplacians_cnn}), homología persistente, complejos simpliciales) proporcionan un marco matemáticamente riguroso para capturar relaciones de orden superior. No obstante, las arquitecturas que operan sobre estas estructuras se encuentran en un estado de madurez temprano en comparación con las arquitecturas sobre grafos. El diseño de operaciones de convolución y atención eficientes sobre complejos celulares generales, que preserven las propiedades algebraicas de la topología subyacente sin incurrir en costes computacionales prohibitivos, permanece como un problema abierto relevante. El artículo de posición de \citet{papamarkou2024tdlposition}, firmado por buena parte de la comunidad, enumera los problemas abiertos del \gls{TDL} en estos términos: una teoría de la expresividad sobre complejos análoga a la jerarquía WL para grafos, la escalabilidad del paso de mensajes sobre celdas de dimensión alta, la falta de benchmarks en los que la estructura topológica sea demostrablemente necesaria, y la conexión con los sheaves celulares y otras estructuras algebraicas; el survey de \citet{papillon2023tdlarchitectures} organiza las arquitecturas existentes bajo una única plantilla de paso de mensajes sobre hipergrafos y complejos simpliciales, celulares y combinatorios, lo que hace posible una comparación sistemática entre ellas.


\section{Hacia Dónde Va el Campo}

Las direcciones que siguen se fundamentan directamente en los resultados y limitaciones identificados arriba. Se enuncian como perspectivas, no como resultados: el libro no afirma sobre ellas nada más allá de lo que establece la literatura citada.

\subsection{Arquitecturas Híbridas CNN-Atención}

El teorema de correspondencia (\autoref{thm:correspondencia_cnn_attention}) y la conjetura de convergencia (\autoref{conj:convergencia_cnn_attention}) entre convoluciones y atención, junto con el análisis del espectro rigidez-adaptabilidad (\autoref{rem:tradeoff_rigidez_adaptabilidad}), sugieren que una dirección natural es el diseño de arquitecturas que combinen ambos paradigmas de forma controlada. El marco teórico establecido en el \autoref{ch:attention} proporciona las herramientas formales para interpolar entre la rigidez de las convoluciones equivariantes y la adaptabilidad de la atención, potencialmente combinando la eficiencia en el uso de datos de las primeras con la expresividad de la segunda.

Dos patrones ilustran la interpolación. El primero empareja un núcleo convolucional, que procesa características locales con garantías geométricas, con cabezas de atención que capturan dependencias globales adaptativas, y deja que una compuerta aprendida decida, por capa o por token, cuánto usar de cada uno; las arquitecturas híbridas de visión del \autoref{ch:attention} son instancias tempranas. El segundo sustituye el softmax de la atención por un decaimiento fijo con la distancia en el dominio, como hacen las redes retentivas con la posición en la secuencia, pero con la distancia geodésica $d_\mathcal{M}$ del \autoref{ch:gdl} en lugar del desplazamiento posicional, de modo que el decaimiento respeta la geometría y no un orden arbitrario. Cuál de los dos preserva mejor las garantías del \autoref{ch:cnn} conservando la flexibilidad del \autoref{ch:attention} es una cuestión empírica que la teoría de este libro enmarca pero no resuelve.

\subsection{Dominios Geométricos Complejos}

El marco de las 5~Gs (\autoref{subsec:cinco_gs}) admite extensiones naturales a espacios producto que combinen geometrías de distinta curvatura (hiperbólica, esférica y euclídea) en un mismo modelo. Estas estructuras surgen de forma natural en datos con componentes jerárquicas y cíclicas simultáneas. En la dirección topológica, los resultados sobre complejos simpliciales y sheaves celulares del \autoref{ch:cnn} proporcionan una base para el desarrollo de \textit{sheaf neural networks}, que generalizarían las redes de grafos al incorporar la estructura algebraica de los haces sobre complejos celulares.

Una tercera extensión concierne a las propias simetrías y no al dominio. El marco de las 5~Gs asume un grupo de simetría global, mientras que la mayoría de los dominios reales solo ofrecen simetrías parciales: isomorfismos locales entre vecindades de un grafo, o isometrías radiales entre parches de una variedad. Los grupoides (\autoref{subsubsec:grupoides}) proporcionan la estructura algebraica para estas simetrías, y trabajos recientes establecen convoluciones equivariantes y teoremas de aproximación universal sobre grupoides de Lie \citep{astwood2026liegroupoid, maruyama2025categorical}, mientras que el análisis de las simetrías de fibración \citep{velarde2026fibration} muestra que las simetrías locales no automórficas de un grafo ya restringen qué funciones puede calcular una red de paso de mensajes. Si estos formalismos producen arquitecturas que superen en la práctica a las Natural Graph Networks \citep{dehaan2020natural} y a las redes gauge equivariantes sigue siendo una cuestión abierta: la teoría está disponible, pero aún no existen benchmarks.

\subsection{Modelos Generativos sobre Variedades y Grupos de Lie}\label{subsec:modelos_generativos_variedades}

El material sobre variedades y grupos de Lie del \autoref{ch:gdl} (la métrica riemanniana (\autoref{def:metrica_riemanniana}), las geodésicas (\autoref{def:geodesica}), el mapa exponencial (\autoref{def:mapa_exponencial_riemanniano}) y el operador de Laplace-Beltrami (\autoref{def:laplace_beltrami})) se ha utilizado en este libro para construir operadores equivariantes discriminativos. Esas mismas herramientas son exactamente lo que se necesita para modelar distribuciones de probabilidad sobre espacios no euclidianos, una dirección que el libro no ha cubierto y que se ha convertido en una de las aplicaciones más activas del \gls{GDL}. El único ingrediente que falta es estocástico: el movimiento browniano sobre una variedad, la difusión cuyo generador es $\tfrac{1}{2}\Delta_\mathcal{M}$. Tres construcciones ilustran esta dirección.

\textbf{Modelado generativo riemanniano basado en el score.} \citet{debortoli2022riemanniansgm} extienden los modelos de difusión basados en el score a una variedad riemanniana compacta $(\mathcal{M}, g)$ sustituyendo el proceso euclidiano de adición de ruido por el movimiento browniano riemanniano $\mathrm{d}\mathbf{X}_t = \mathrm{d}\mathbf{B}^\mathcal{M}_t$, cuya densidad de transición $p_t(\cdot \mid \mathbf{x}_0)$ es el núcleo del calor (\textit{heat kernel}) del operador de Laplace-Beltrami, $\partial_t p_t = \tfrac{1}{2}\Delta_\mathcal{M} p_t$, y que converge a la distribución uniforme respecto del volumen riemanniano. Una red $s_\theta(t, \mathbf{x}) \in T_\mathbf{x}\mathcal{M}$ se entrena mediante \textit{score matching} para aproximar el score riemanniano $\nabla_\mathbf{x} \log p_t(\mathbf{x})$, un campo vectorial tangente, y las muestras se generan integrando la SDE invertida en el tiempo $\mathrm{d}\mathbf{Y}_t = s_\theta(T - t, \mathbf{Y}_t)\,\mathrm{d}t + \mathrm{d}\mathbf{B}^\mathcal{M}_t$, discretizada como un paseo aleatorio geodésico $\mathbf{Y}_{k+1} = \exp_{\mathbf{Y}_k}\!\big(\gamma\, s_\theta(T - t_k, \mathbf{Y}_k) + \sqrt{\gamma}\,\boldsymbol{\xi}_k\big)$ con $\boldsymbol{\xi}_k$ gaussiano en $T_{\mathbf{Y}_k}\mathcal{M}$. En espacios simétricos como la esfera, el núcleo del calor es una serie explícita en las autofunciones de $\Delta_\mathcal{M}$ (los armónicos esféricos), lo que hace computable el objetivo de entrenamiento; en variedades generales debe aproximarse, y este es el principal coste práctico del método.

\textbf{Flow matching riemanniano.} \citet{chen2024flowmatchinggeometries} evitan por completo la simulación estocástica. Dada una distribución fuente $p_0$ sobre $\mathcal{M}$ (típicamente uniforme) y la distribución de los datos $p_1$, se muestrean $\mathbf{x}_0 \sim p_0$, $\mathbf{x}_1 \sim p_1$ y se conectan mediante el interpolante geodésico
\begin{equation}
    \mathbf{x}_t = \exp_{\mathbf{x}_0}\!\big(t \log_{\mathbf{x}_0}(\mathbf{x}_1)\big), \qquad t \in [0, 1],
\end{equation}
cuya velocidad $\dot{\mathbf{x}}_t = \log_{\mathbf{x}_t}(\mathbf{x}_1)/(1 - t)$ es el vector tangente que apunta desde $\mathbf{x}_t$ hacia el dato. Un campo vectorial neuronal $v_\theta(t, \cdot)$ tangente a $\mathcal{M}$ se entrena por regresión sobre esta velocidad,
\begin{equation}\label{eq:riemannian_flow_matching}
    \mathcal{L}_{\text{RFM}}(\theta) = \mathbb{E}_{t \sim \mathcal{U}[0,1],\; \mathbf{x}_0 \sim p_0,\; \mathbf{x}_1 \sim p_1}\, \big\| v_\theta(t, \mathbf{x}_t) - \dot{\mathbf{x}}_t \big\|^2_{g(\mathbf{x}_t)},
\end{equation}
donde $\|\cdot\|_{g}$ es la norma inducida por la métrica riemanniana. Aunque el objetivo de regresión es condicional al par $(\mathbf{x}_0, \mathbf{x}_1)$, el minimizador es el campo vectorial marginal que transporta $p_0$ a $p_1$, y el muestreo se reduce a integrar la ODE $\dot{\mathbf{x}} = v_\theta(t, \mathbf{x})$ sobre la variedad con pasos del mapa exponencial. El entrenamiento no requiere simulación siempre que $\exp$ y $\log$ estén disponibles en forma cerrada: la esfera (\autoref{ex:mapa_exp_esfera}), el toro plano, el espacio hiperbólico y $SO(3)$ mediante la fórmula de Rodrigues (\autoref{ex:mapa_exp_so3}). En geometrías generales (mallas con frontera, por ejemplo), la distancia geodésica del interpolante se sustituye por una premétrica calculada a partir de una descomposición espectral del operador de Laplace-Beltrami, al precio de aproximar el campo vectorial objetivo.

\textbf{Difusión en $SE(3)$ para esqueletos de proteínas.} Un esqueleto de $N$ residuos se representa, siguiendo la convención de AlphaFold, mediante $N$ marcos rígidos $\mathbf{T}_i = (\mathbf{R}_i, \mathbf{x}_i) \in SE(3)$, donde la rotación orienta la tríada local de átomos del esqueleto y la traslación sitúa el átomo $\text{C}_\alpha$; un esqueleto es así un punto del grupo de Lie $SE(3)^N$. \citet{yim2023se3diffusion} construyen la difusión directamente a partir de la estructura de grupo $SE(3) = SO(3) \ltimes \mathbb{R}^3$: con una métrica invariante por la izquierda adecuada, el movimiento browniano en $SE(3)$ se desacopla en un movimiento browniano euclidiano sobre las traslaciones (tras eliminar el centro de masas, para que exista una distribución estacionaria) y un movimiento browniano en $SO(3)$ con su métrica bi-invariante, cuya densidad de transición es la gaussiana isótropa en $SO(3)$.

\begin{definition}[Gaussiana isótropa en $SO(3)$]\label{def:igso3}
La distribución gaussiana isótropa $\mathcal{IG}_{SO(3)}(\mathbf{R}; \mathbf{R}_0, \sigma^2)$ es la densidad de transición, respecto de la medida de Haar normalizada, del movimiento browniano en $SO(3)$ con la métrica bi-invariante, iniciado en $\mathbf{R}_0$ y ejecutado durante un tiempo $\sigma^2$. Depende únicamente del ángulo de rotación $\omega \in [0, \pi]$ de la rotación relativa $\mathbf{R}_0^\top \mathbf{R}$, dado por $\cos\omega = (\operatorname{tr}(\mathbf{R}_0^\top \mathbf{R}) - 1)/2$, a través de
\begin{equation}\label{eq:igso3_densidad}
    f(\omega; \sigma^2) = \sum_{\ell = 0}^{\infty} (2\ell + 1)\, e^{-\ell(\ell+1)\sigma^2/2}\, \frac{\sin\big((\ell + \tfrac{1}{2})\omega\big)}{\sin(\omega/2)}.
\end{equation}
\end{definition}

La serie de \eqref{eq:igso3_densidad} es la expansión de Peter-Weyl (\autoref{thm:peter_weyl}) de una función de clase en $SO(3)$: $\chi_\ell(\omega) = \sin((\ell + \tfrac{1}{2})\omega)/\sin(\omega/2)$ es el carácter de la representación irreducible de dimensión $2\ell + 1$, y $-\ell(\ell+1)$ es el autovalor del operador de Laplace-Beltrami sobre la componente isotípica correspondiente. El núcleo del calor de cualquier grupo de Lie compacto tiene la misma forma, $\sum_\rho d_\rho\, e^{-\lambda_\rho t}\, \chi_\rho$, razón por la cual el movimiento browniano en tal grupo es computable a partir únicamente de su tabla de caracteres. El score de \eqref{eq:igso3_densidad} se obtiene derivando la serie respecto de $\omega$, la red de score produce un elemento del álgebra de Lie $\mathfrak{so}(3)$ por residuo, devuelto al grupo mediante el mapa exponencial, y la red se hace $SE(3)$-equivariante para que la distribución generada sobre esqueletos sea $SE(3)$-invariante: un campo vectorial equivariante iniciado desde una distribución invariante la transporta a otra invariante.

En este último punto reside el problema abierto. Incorporar la equivarianza al modelo generativo, como arriba, proporciona invarianza exacta de la distribución de las muestras y eficiencia en datos; AlphaFold~3 \citep{abramson2024alphafold3}, en cambio, descartó el módulo de estructura $SE(3)$-equivariante de su predecesor en favor de un módulo de difusión que opera sobre coordenadas atómicas brutas sin marcos ni procesamiento equivariante, y recupera la simetría solo de forma aproximada mediante rotaciones y traslaciones aleatorias de los datos de entrenamiento. Si el aumento de datos hace innecesarios los modelos generativos equivariantes restringidos, y a qué escala de datos y de modelo, es la contrapartida generativa de la tensión rigidez-adaptabilidad discutida en el \autoref{ch:attention} (\autoref{rem:tradeoff_rigidez_adaptabilidad}) y del dilema equivarianza-escalabilidad analizado en la \autoref{subsubsec:dilema_equivarianza_escalabilidad}.

\subsection{Análisis Armónico Generalizado}

El papel central del análisis armónico en este libro (desde las convoluciones espectrales hasta los armónicos esféricos y la descomposición de Chebyshev del \autoref{ch:cnn}) sugiere la exploración de extensiones a clases más amplias de grupos. El teorema de Peter-Weyl (\autoref{thm:peter_weyl}), piedra angular de las convoluciones de grupo, requiere compacidad del grupo; la extensión del marco a grupos no compactos (como el grupo de Lorentz o los grupos afines) demanda herramientas de análisis armónico no conmutativo que van más allá de las desarrolladas en este libro. Un avance en esta dirección ampliaría el ámbito de aplicabilidad de las \glspl{CNN} equivariantes a simetrías relevantes en física relativista y procesamiento de señales generales.


\subsection{Hacia una Teoría Unificada}\label{subsec:hacia_teoria_unificada}

Dos síntesis recientes cartografían el terreno que una teoría unificada del \gls{GDL} tendría que reconciliar: \citet{saezdeocarizborde2025mathfoundations} organizan el campo desde el lado del modelo, en torno a las estructuras matemáticas (grupos de simetría, variedades, grafos) impuestas a la arquitectura, mientras que \citet{papillon2024beyondeuclid} lo organizan desde el lado de los datos, clasificando los métodos por la estructura geométrica, topológica y algebraica de sus entradas; una teoría unificada debería explicar por qué ambas taxonomías coinciden allí donde lo hacen. Tres preguntas, abiertas pero no respondidas en este libro, señalan el camino.

La primera concierne a los límites de la expresividad. La conjetura de convergencia \gls{CNN}-atención (\autoref{conj:convergencia_cnn_attention}) afirma que ambos paradigmas convergen al mismo conjunto límite de operadores. Si es así, cabría esperar que también converjan sus costes de aproximación: para un objetivo continuo sobre una variedad compacta, el tamaño de la menor red convolucional y el de la menor red de atención que alcanzan una precisión dada deberían hacerse comparables a medida que el requisito de precisión se endurece. El teorema de universalidad equivariante (\autoref{theorem:equivariant_universality_cnns}) garantiza que ambos tamaños son finitos; nada en la teoría actual los relaciona, y falta una comparación cuantitativa de las tasas de aproximación entre paradigmas.

La segunda es una teoría de la información propiamente geométrica. La cota de complejidad muestral de la \autoref{sec:complejidad_muestral} muestra que lo que importa es la dimensión intrínseca de la variedad de datos, no la ambiente, y la cota de \gls{oversquashing} de \citet{digiovanni2023oversquashing} muestra que lo que limita el flujo de información en una red de grafos es la resistencia efectiva del grafo. Ambos resultados ponderan una cantidad estadística por una geométrica. Una teoría que hiciera esto de forma sistemática, midiendo la información que porta una representación respecto del volumen riemanniano del dominio y no respecto de la medida de Lebesgue, daría una forma principiada de comparar arquitecturas entre dominios y de decidir cuánta estructura geométrica requiere realmente una tarea. Las piezas existen en la geometría de la información; su aplicación al diseño de redes equivariantes, no.

La tercera es el descubrimiento automático de la simetría. A lo largo del libro el grupo ha venido dado: traslaciones para imágenes, $SE(3)$ para moléculas, permutaciones para conjuntos. El frame averaging y la canonicalización (\autoref{subsec:canonicalizacion_frames}) ya desacoplan la equivarianza de la arquitectura, y los métodos de descubrimiento de simetrías como LieGAN \citep{yang2023liegan} aprenden los generadores de un álgebra de Lie de simetrías a partir de los datos. Juntos, esbozan una cadena que va del conjunto de datos a su grupo de simetría y de este a la arquitectura equivariante correspondiente sin intervención humana. Si el grupo así descubierto es el que habría elegido un matemático, y si la arquitectura resultante hereda las garantías demostradas en este libro para grupos conocidos, son preguntas abiertas.

Estas preguntas comparten una misma forma: para un problema con una estructura geométrica conocida y un grado de adaptabilidad requerido, piden la arquitectura que respeta la estructura hasta donde esta se cumple y se adapta donde no. El espectro rigidez-adaptabilidad del \autoref{ch:attention} es el eje a lo largo del cual se situaría tal arquitectura; la teoría unificada sería la regla que la sitúa.


\section{Reflexión Final}

El desarrollo presentado en este libro muestra que la geometría no es simplemente una herramienta para mejorar el rendimiento de los modelos de aprendizaje profundo, sino un principio organizador que permite explicar por qué determinadas arquitecturas funcionan para determinados tipos de datos. Desde esta perspectiva teórica profunda, la distinción entre lo ``convolucional'' y lo ``atencional'' se vuelve cada vez menos relevante: ambos paradigmas son puntos distintos de un mismo continuo de operadores geométricos, y el futuro del campo reside menos en su competencia que en su síntesis principiada. La síntesis del \autoref{ch:gdl} (\autoref{sec:sintesis_marco_unificador}) ha mostrado que las diferencias entre \glspl{ANN}, \glspl{CNN}, \glspl{GNN} y \glspl{Transformer} reflejan elecciones sobre tres aspectos fundamentales: el dominio geométrico, el grupo de simetría y la estructura de conectividad.

El marco matemático que sostiene esta unificación (teoría de Lie, geometría riemanniana, teoría gauge, topología algebraica) posee un poder que trasciende la mera descripción retrospectiva de las arquitecturas existentes. Los resultados de universalidad equivariante y los teoremas de convergencia entre paradigmas proporcionan herramientas con capacidad predictiva: permiten anticipar qué restricciones geométricas son necesarias para un dominio dado y qué garantías expresivas ofrecerá la arquitectura resultante. En este sentido, el \gls{GDL} constituye no solo una taxonomía de las arquitecturas actuales, sino un marco teórico para el diseño principiado de las futuras.
